\chapter{Chapter 7 Refresher: Efficiency, Scaling, and Mixture-of-Experts --- Systems Arithmetic and Sparse Routing}

\section*{Setting the Scene}
Once probabilities and networks are defined, training becomes physics on silicon: bytes move, joules dissipate, routers stall. Preface Step~6 (variance of sums) quietly lurks behind contention noise.

\paragraph{Step A --- Roofline worldview.}
\textbf{Problem.} MFLOPS dashboards lie when kernels stream weights constantly.
\textbf{Insight.} Compare arithmetic intensity (FLOPs/byte) against bandwidth-induced slopes---whichever bound bites first wins.
\textbf{Lineage.} Williams/Waters/Patterson roofline model (2009) repurposed across ML hardware discourse.
\textbf{Mental model.} If intensity $<$ ridge crossover, buying more TFLOPs wastes money---buy bandwidth or locality instead.

\paragraph{Step B --- Sparsity swaps villains.}
\textbf{Problem.} Dense matrices saturate compute; sparse routing saves FLOPs but reroutes traffic.
\textbf{Insight.} MoE trades matmuls for all-to-all collectives---new bottleneck on wires (Preface Step~6 covariance spikes across devices).
\textbf{Lineage.} Shazeer's Switch Transformer et seq.
\textbf{Mental model.} Plot effective \(\lambda/\mu\) including dispatch latency before celebrating parameter counts.

\paragraph{Step C --- Capacity factor as queue abstraction.}
\textbf{Problem.} Experts overload like servers under burst arrivals.
\textbf{Insight.} Capacity factor overprovisions slots analogous to multi-server queue stability margins---drops mirror admission control.
\textbf{Lineage.} Queueing tradition (Kendall notation) meets routing papers.
\textbf{Mental model.} Tail latency dominates UX---average FLOPs/token misleads.

Systems graphs reinterpret probability routines as throughput limits; migrating bottlenecks obey no moral ranking beyond measurement.

\section*{Notation Introduced in This Chapter}
\begin{itemize}[leftmargin=1.5em]
  \item $I$: arithmetic intensity (FLOPs per byte).
  \item $k$: number of selected experts per token in top-$k$ routing.
  \item $N_{\text{active}}$: active parameters per token.
  \item $N_{\text{total}}$: total resident parameters.
  \item CF: capacity factor (expert slot overprovision multiplier).
\end{itemize}

\section*{What You Need To Recall Precisely}
\begin{itemize}[leftmargin=1.5em]
  \item \textbf{Bottleneck triad:} compute, memory bandwidth, and interconnect communication.
  \item \textbf{MoE routing:} compute gate logits per token (router network); softmax yields mixture weights; keep top-$k$ experts, zero or mask others, renormalise weights on the surviving experts so each token still mixes convexly over experts it dispatches to; aggregate expert outputs with those weights. Routing induces sparse activation patterns separate from parameter residency.
  \item \textbf{Capacity factor:} finite per-expert token slots induce drop/reroute decisions.
  \item \textbf{Arithmetic intensity:} $\text{FLOPs}/\text{byte}$ predicts memory-bound vs compute-bound behavior.
  \item \textbf{Active vs total parameters:} sparse models can have large total memory footprint while using small active compute per token.
\end{itemize}

\section*{How These Results Build on Each Other}
\begin{enumerate}[leftmargin=1.5em]
  \item Compute workload demand (FLOPs and bytes/token).
  \item Compare with device and interconnect limits.
  \item Identify bottleneck (compute, memory, communication).
  \item Apply sparse routing and reassess new bottleneck.
  \item Validate with throughput/latency measurements.
\end{enumerate}

\section*{Reminder Glossary (Term by Term)}
\begin{itemize}[leftmargin=1.5em]
  \item \textbf{Arithmetic intensity:} FLOPs per byte transferred.
  \item \textbf{Dispatch/combine:} route tokens to experts and gather outputs.
  \item \textbf{Capacity factor:} overprovisioning multiplier for expert token slots.
  \item \textbf{Token drop:} overflow handling when expert capacity exceeded.
  \item \textbf{All-to-all:} communication primitive central to MoE scaling.
\end{itemize}

\section*{Worked Mini Example (Intensity Intuition)}
Assume a kernel performs $2\times10^{12}$ FLOPs and moves $8\times10^{11}$ bytes.
\[
I=\frac{2\times10^{12}}{8\times10^{11}}=2.5\ \text{FLOPs/byte}.
\]
If hardware requires much higher intensity for compute saturation, this kernel is memory-bound.

\section*{Worked Example (MoE Bottleneck Migration)}
Assume dense model requires 100 units compute and 40 units communication per batch.  
After MoE sparsification, active compute drops to 45, but dispatch/combine communication rises to 80.  
System shifts from compute-bound to communication-bound.  
Interpretation: sparse compute improved one axis but exposed another, so optimization must target network scheduling or expert placement.

\section*{Intuition Checkpoints}
\begin{itemize}[leftmargin=1.5em]
  \item Sparse compute does not imply sparse communication.
  \item MoE can reduce compute/token while increasing routing complexity.
  \item High model parameter count can coexist with low active compute and high memory residency.
\end{itemize}

\section*{Further Refresh}
\begin{itemize}[leftmargin=1.5em]
  \item Optimization and memory accounting: see Chapter 2 refresher.
  \item Attention complexity and KV behavior: see Chapter 5 refresher.
  \item Decoding/token-rate implications: see Chapter 6 refresher.
  \item \textbf{Shazeer et al. / Switch Transformer papers} --- sparse routing mechanics and load balancing.
  \item \textbf{Harchol-Balter performance modeling} --- bottleneck and throughput reasoning.
\end{itemize}
